% Part: lambda-calculus
% Chapter: introduction
% Section: minimization

\documentclass[../../../include/open-logic-section]{subfiles}

\begin{document}

\olfileid[id]{lam}{int}{min}
\olsection{Fungsi yang \usetoken{S}{lambda definable} Tertutup terhadap Minimisasi}

\begin{lem}
Andaikan $f(x,y)$ bersifat !!{lambda definable}. Misalkan $g$ didefinisikan
oleh
\[
g(x) \simeq \umin{y}{f(x,y)}.
\]
Maka $g$ bersifat !!{lambda definable}.
\end{lem}

\begin{proof}
Gagasannya kira-kira sebagai berikut. Jika $x$ diberikan, kita akan
menggunakan suku lambda titik-tetap $Y$ untuk mendefinisikan fungsi $h_x(n)$
yang mencari suatu~$y$ mulai dari~$n$; kemudian $g(x)$ hanyalah $h_x(0)$.
Fungsi~$h_x$ dapat dinyatakan sebagai solusi persamaan titik tetap:
\[
h_x(n) \simeq
\begin{cases}
n & \text{jika $f(x,n) = 0$} \\
h_x(n+1) & \text{dalam kasus lainnya.}
\end{cases}
\]

Berikut perinciannya. Karena $f$ bersifat !!{lambda definable}, fungsi itu
!!{lambda defined} oleh suatu suku~$F$. Ingat bahwa kita juga mempunyai suku
lambda~$D$ sedemikian sehingga $D(M,N,\bar{0}) \red M$ dan
$D(M,N,\bar{1}) \red N$. Dengan menetapkan $x$ untuk sementara, untuk
!!{lambda define} $h_x$ kita ingin menemukan suatu suku~$H$ (yang bergantung
pada~$x$) yang memenuhi
\[
H(\num{n}) \equiv D(\num{n}, H(S(\num{n})), F(x, \num{n})).
\]
Kita dapat melakukannya dengan menggunakan suku titik-tetap~$Y$. Pertama,
misalkan $U$ adalah suku
\[
\lambd[h][\lambd[z][D(z,(h(Sz)),F(x,z))]],
\]
kemudian misalkan $H$ adalah suku~$YU$. Perhatikan bahwa satu-satunya variabel
bebas dalam~$H$ adalah~$x$. Mari kita tunjukkan bahwa $H$ memenuhi persamaan
di atas.

Berdasarkan definisi~$Y$, kita mempunyai
\[
H = YU \equiv U(YU) = U(H).
\]
Khususnya, untuk setiap bilangan asli~$n$, kita mempunyai
\begin{align*}
H(\num{n}) & \equiv U(H, \num{n}) \\
& \red D(\num{n}, H(S(\num{n})), F(x, \num{n})),
\end{align*}
seperti yang diperlukan. Perhatikan bahwa jika numeral $\num{m}$
disubstitusikan sebagai pengganti~$x$ pada baris terakhir, ekspresi itu
tereduksi menjadi $\num{n}$ jika $F(\num{m},\num{n})$ tereduksi menjadi
$\num{0}$, dan tereduksi menjadi $H(S(\num{n}))$ jika
$F(\num{m},\num{n})$ tereduksi menjadi numeral lainnya.

Untuk menyelesaikan pembuktian, misalkan $G$ adalah
$\lambd[x][H(\num 0)]$. Maka $G$ !!{lambda define}s~$g$; dengan kata lain,
untuk setiap~$m$, $G(\num m)$ tereduksi menjadi~$\overline{g(m)}$ jika $g(m)$
terdefinisi, dan tidak memiliki bentuk normal dalam kasus lainnya.
\end{proof}

\end{document}
